\documentclass[11pt,a4paper]{article}

\usepackage[a4paper,margin=1in]{geometry}
\usepackage{graphicx}
\usepackage{booktabs}
\usepackage{float}
\usepackage{array}
\usepackage{hyperref}
\usepackage{enumitem}

\hypersetup{
  colorlinks=true,
  linkcolor=black,
  urlcolor=blue
}

\title{PDE4444 Visual Quality Inspection Report\\Cup Surface Defect Classification}
\author{Student Name: \underline{\hspace{5cm}}}
\date{\today}

\begin{document}
\maketitle

\section*{1. Problem Definition and Ground Truth}
This project implements a visual quality inspection system for cups using a single camera image and a binary machine-learning decision. The selected inspection rule is \textbf{visible surface defect present or absent}. This keeps the submission aligned to one clear engineering rule and one final decision per image.

\begin{itemize}[leftmargin=*]
  \item \textbf{PASS} = cup is \texttt{non\_defective}
  \item \textbf{FAIL} = cup is \texttt{defective}
\end{itemize}

Ground truth labels were assigned by manual inspection under a fixed capture setup. The intended deployment flow is: object placed in inspection area $\rightarrow$ image captured $\rightarrow$ classifier predicts \texttt{defective}/\texttt{non\_defective} $\rightarrow$ system maps output to \textbf{FAIL}/\textbf{PASS}.

\begin{figure}[H]
\centering
\fbox{\parbox{0.88\linewidth}{\centering Add final inspection-area or stream screenshot here before submission.}}
\caption{Optional figure slot for the fixed inspection setup or live stream interface.}
\end{figure}

\section*{2. Model Design and Training}
\subsection*{Dataset}
The dataset uses two raw classes: \texttt{defective} and \texttt{non\_defective}. Augmentation was applied offline and cleaned before splitting in order to reduce exact and near-duplicate samples.

\begin{table}[H]
\centering
\begin{tabular}{lrr}
\toprule
\textbf{Stage / Split} & \textbf{Defective} & \textbf{Non-defective} \\
\midrule
Raw captures & 67 & 7 \\
Cleaned pre-split dataset & 340 & 192 \\
Train split & 241 & 139 \\
Validation split & 46 & 29 \\
Test split & 53 & 24 \\
\bottomrule
\end{tabular}
\caption{Dataset counts after augmentation, presplit cleanup, and source-group splitting.}
\end{table}

The split ratio was 70/15/15 by source group in order to reduce leakage from augmented siblings. This matters because multiple augmented samples were derived from the same original source image.

\subsection*{Models}
\begin{itemize}[leftmargin=*]
  \item \textbf{Main model}: YOLO26 image classification with pretrained weights (\texttt{yolo26n-cls.pt})
  \item \textbf{Baseline model}: HOG features with Logistic Regression
\end{itemize}

Training settings for the main model were:
\begin{itemize}[leftmargin=*]
  \item image size: 224
  \item epochs: 80
  \item batch size: 16
  \item patience: 15
  \item pretrained initialization enabled
\end{itemize}

\subsection*{Training loop pseudo-code}
\begin{verbatim}
load train and validation datasets
initialize pretrained classification model
for epoch in range(num_epochs):
    model.train()
    for images, labels in train_loader:
        predictions = model(images)
        loss = criterion(predictions, labels)
        optimizer.zero_grad()
        loss.backward()
        optimizer.step()

    model.eval()
    with no_grad():
        run validation predictions
        compute validation metrics
save best model weights
\end{verbatim}

\section*{3. Quantitative Results}
\subsection*{Validation and baseline summary}
The YOLO26 model reached a best validation top-1 accuracy of \textbf{96.0\%} at epoch 10 in the first validation run. The HOG + Logistic Regression baseline reached \textbf{79.0\%} accuracy on the validation split, so the CNN classifier was clearly the stronger approach.

\begin{table}[H]
\centering
\begin{tabular}{lcccc}
\toprule
\textbf{Model} & \textbf{Split} & \textbf{Accuracy} & \textbf{Macro Precision} & \textbf{Macro F1} \\
\midrule
HOG + Logistic Regression & Validation & 0.790 & 0.800 & 0.750 \\
YOLO26n classifier & Test & 0.896 & 0.872 & 0.884 \\
\bottomrule
\end{tabular}
\caption{Baseline sanity check versus final YOLO classifier.}
\end{table}

\subsection*{Final test-set results}
The saved best model (\texttt{models/PDE4444\_T1\_best.pt}) was evaluated on the held-out test split using the project evaluation script. The resulting report was:

\begin{itemize}[leftmargin=*]
  \item Accuracy: 0.8961
  \item Macro precision: 0.8724
  \item Macro recall: 0.9017
  \item Macro F1-score: 0.8839
  \item Weighted F1-score: 0.8981
\end{itemize}

Per-class test performance:
\begin{itemize}[leftmargin=*]
  \item \texttt{non\_defective} (PASS): precision 0.7857, recall 0.9167, F1 0.8462
  \item \texttt{defective} (FAIL): precision 0.9592, recall 0.8868, F1 0.9216
\end{itemize}

Confusion matrix (rows=true, cols=pred):
\[
\begin{bmatrix}
22 & 2 \\
6 & 47
\end{bmatrix}
\]

\begin{figure}[H]
\centering
\includegraphics[width=0.8\linewidth]{../1st_validation_training/confusion_matrix.png}
\caption{Confusion matrix from the first validation training run.}
\end{figure}

\begin{figure}[H]
\centering
\includegraphics[width=0.85\linewidth]{../1st_validation_training/results.png}
\caption{Training and validation curves from the YOLO26 classification run.}
\end{figure}

\section*{4. Failure Cases and Limitations}
The model is usable for the PASS/FAIL requirement, but the current dataset still has important limitations.

\begin{itemize}[leftmargin=*]
  \item The biggest limitation is the low diversity of real PASS examples: only 7 raw \texttt{non\_defective} source groups were available.
  \item The test confusion matrix shows 6 defective cups predicted as PASS, which is the more critical failure mode for quality inspection.
  \item Likely causes include glare, low-contrast defects, unusual cup pose, and insufficient real-world variation in the non-defective class.
  \item Augmentation improved volume, but synthetic diversity does not fully replace additional real captures.
\end{itemize}

\begin{figure}[H]
\centering
\includegraphics[width=0.8\linewidth]{../1st_validation_training/val_batch0_pred.jpg}
\caption{Example validation predictions from the training run.}
\end{figure}

\begin{figure}[H]
\centering
\fbox{\parbox{0.88\linewidth}{\centering Add one or two final failure-case screenshots here before submission.}}
\caption{Optional figure slot for selected false PASS / false FAIL examples.}
\end{figure}

\subsection*{Final engineering output}
The report and demo should always present the model output as an engineering decision:
\begin{itemize}[leftmargin=*]
  \item \texttt{non\_defective} $\rightarrow$ \textbf{PASS}
  \item \texttt{defective} $\rightarrow$ \textbf{FAIL}
\end{itemize}

For the final demonstration, at least 5 genuinely unseen samples should be kept outside the training pipeline and evaluated with the same PASS/FAIL mapping.

\end{document}
